\documentclass[aps,pra,reprint,superscriptaddress,longbibliography,nofootinbib,floatfix,10pt]{revtex4-2}
\usepackage[T1]{fontenc}
\usepackage{lmodern,amsmath,amssymb,amsthm,mathtools,bm,graphicx,booktabs,longtable,tikz,placeins,needspace,microtype}
\usetikzlibrary{calc,arrows.meta}
\usepackage[colorlinks=true,linkcolor=blue,citecolor=blue,urlcolor=blue,breaklinks=true]{hyperref}
\usepackage{orcidlink}

\newtheorem{theorem}{Theorem}
\newtheorem{lemma}[theorem]{Lemma}
\newtheorem{proposition}[theorem]{Proposition}
\newtheorem{corollary}[theorem]{Corollary}
\theoremstyle{definition}\newtheorem{definition}[theorem]{Definition}
\theoremstyle{remark}\newtheorem{remark}[theorem]{Remark}

\newcommand{\C}{\mathbb C}\newcommand{\R}{\mathbb R}\newcommand{\Q}{\mathbb Q}
\newcommand{\F}{\mathbb F}\newcommand{\Z}{\mathbb Z}\newcommand{\PP}{\mathbb P}
\newcommand{\Sc}{\mathcal S}\newcommand{\Cc}{\mathcal C}
\newcommand{\YO}{\mathrm{YO}}\newcommand{\ip}[2]{\langle #1,#2\rangle}
\DeclareMathOperator{\spanop}{span}\DeclareMathOperator{\rank}{rank}
\DeclareMathOperator{\diag}{diag}\DeclareMathOperator{\Tr}{Tr}\DeclareMathOperator{\tr}{tr}
\DeclareMathOperator{\supp}{supp}\DeclareMathOperator{\conv}{conv}

\hypersetup{
  pdftitle={Context-Color Matrices and Latin-Square Structure in Quantum Measurements},
  pdfauthor={Karl Svozil}
}

\begin{document}

\title{Context--Color Matrices and Latin-Square Structure in Quantum Measurements}

\author{NN}
%\author{Karl Svozil\,\orcidlink{0000-0001-6554-2802}}
%\email{karl.svozil@tuwien.ac.at}
%\affiliation{Institute for Theoretical Physics, TU Wien, Wiedner Hauptstra{\ss}e 8-10/136, A-1040 Vienna, Austria}

\date{\today}



\begin{abstract}
We study common-spectrum spectral assignments for intertwined projective quantum measurements. All worked configurations have faithful orthogonal ray representations in the stated Hilbert space, with coordinates supplied here or in the companion inventory. For $n$ contexts in dimension $d$, an $n\times d$ context--color matrix places the ray labels of each basis in one row and aligns shared projectors in their assigned spectral columns. Its columns are extendable two-valued states at the exact palette; wider palettes require blanks. Real three-dimensional examples exhibit nonextendable states, forced scalar identities between distinct rays, and Born probabilities outside the convex hull of color-class states. Faithful four-dimensional extensions and the realized Yu--Oh, complex Hesse, and real $65$-ray configurations illustrate spectral obstructions with surviving ordinary states. Latin-square structure is used only where it organizes actual measurements or their color signatures. An explicit $49$-ray configuration in $\mathbb C^7$ realizes a cyclic Latin net as $21$ projective measurements; all pairwise orthogonalities and nonorthogonalities are certified exactly.
\end{abstract}
\maketitle



\section{Quantum measurements and common-spectrum assignments}\label{sec:intro}

A complete rank-one projective measurement in $\mathbb F^d$, $\mathbb F\in\{\mathbb R,\mathbb C\}$, is an orthonormal basis of rays with projectors $\{P_v:v\in e\}$ satisfying $\sum_{v\in e}P_v=I_d$. Fix distinct real numbers $\Sigma=\{\lambda_0,\ldots,\lambda_{d-1}\}$. Assigning every outcome a different member of $\Sigma$ defines a nondegenerate observable
\begin{equation}\label{eq:context-observable}
A_e=\sum_{v\in e}\lambda_{c(v)}P_v.
\end{equation}
Every operator commuting with this nondegenerate observable is a function of it~\cite{v-neumann-31,halmos-vs}. We study the additional demand that a projector shared by several measurements receives the same spectral label in each. Quantum mechanics permits observables to be specified separately; the common-spectrum assignment is the extra consistency condition under examination~\cite{svozil-2020-c,svozil-2021-chroma,svozil2026contexts}.

Let $H=(V,E)$ list the ray labels and selected complete measurements, with $V=\bigcup E$, $E\ne\varnothing$, $d\ge2$, and $n=|E|$. Its two-section graph $\Gamma(H)$ joins rays occurring together in a context. A \emph{faithful orthogonal representation} (FOR) assigns distinct normalized rays $[q_v]\in\mathbb P(\mathbb F^d)$ such that $q_u\perp q_v$ exactly when $u,v$ are adjacent. Then $P_v=q_vq_v^\dagger$, and the listed contexts are complete projective measurements. Every worked example below has this geometric foundation. The complete coordinates for the companion configurations and their orthogonality checks are supplied in Ref.~\cite{svozil2026contexts,supplement}; the small introductory example, the pentagon, and the seven-dimensional Latin net have coordinates printed here.

A strong $d$-coloring $c:V\to\{0,\ldots,d-1\}$ is injective on every context and gives Eq.~\eqref{eq:context-observable} consistently. We write $\Cc_d(H)$ for these assignments and $\chi(H)$ for the strong chromatic number. A context contains $d$ mutually orthogonal rays, so $\chi(H)\ge d$. The condition $\chi(H)>d$ excludes the common-spectrum labeling; it does not conflict with the existing quantum coordinates.

A two-valued context state is $s:V\to\{0,1\}$ with $\sum_{v\in e}s(v)=1$ for every $e\in E$. Write $\Sc(H)$ for all such states and $\Sc_{\rm ext}(H)$ for those occurring as a color class of some $c\in\Cc_d(H)$. A state extends precisely when the residual graph $\Gamma(H)[V\setminus\operatorname{supp}s]$ admits $d-1$ colors~\cite{svozil2026contexts}. These binary assignments differ from Born probabilities $p_v=\operatorname{Tr}(\rho P_v)$, which satisfy the same context normalization but need not be convex mixtures of color-class states.

Section~\ref{sec:framework} develops the matrix representation and the Latin/signature notation needed for the physical applications. Section~\ref{sec:latin-for} gives an exact realized Latin net for a seven-level quantum system. Sections~\ref{sec:exact-palette} and~\ref{sec:obstructions} treat coordinatized three- and four-dimensional configurations. No uncoordinatized combinatorial construction is used as a worked example.



\section{Context--color matrices for projective measurements}\label{sec:framework}
\subsection{Ray labels as matrix entries}\label{sec:nxd-matrix}

Order the $n$ actual measurement contexts as rows and use the $d$ members of the common spectrum as column labels. Each entry is a ray label, not a vector component. A projector shared by several contexts must appear in the same column throughout. The number $n$ of measurements is independent of the Hilbert-space dimension $d$.
\begin{theorem}[Universal $n\times d$ Context--Color Matrix Criterion]\label{thm:nxd-matrix}
Let $H=(V,E)$ be any finite $d$-uniform hypergraph with $n=|E|$ contexts $E=\{C_1,\ldots,C_n\}$. Then $H$ admits a strong $d$-coloring ($\chi(H)=d$) if and only if the elements of each context $C_i$ can be ordered into an $n\times d$ matrix
\begin{equation}\label{eq:nxd-matrix}
M=\begin{pmatrix}
M_{1,0} & M_{1,1} & \cdots & M_{1,d-1}\\
M_{2,0} & M_{2,1} & \cdots & M_{2,d-1}\\
\vdots  & \vdots  & \ddots & \vdots\\
M_{n,0} & M_{n,1} & \cdots & M_{n,d-1}
\end{pmatrix}\in V^{n\times d},
\end{equation}
satisfying the following two conditions:
\begin{enumerate}
\item \emph{Row Bijectivity (Rainbow Contexts):} For each row $i\in\{1,\ldots,n\}$, the row entries $\{M_{i,0},M_{i,1},\ldots,M_{i,d-1}\}=C_i$ are a permutation of the $d$ distinct vertices of context $C_i$.
\item \emph{Column Alignment (Noncontextual Intertwining Consistency):} Every vertex $v\in V$ appears in \emph{exactly one column} of $M$:
\begin{equation}\label{eq:col-alignment}
M_{i,j}=M_{i',j'}=v \quad\Longrightarrow\quad j=j'.
\end{equation}
\end{enumerate}
When Eqs.~\eqref{eq:nxd-matrix}--\eqref{eq:col-alignment} hold, setting $c(M_{i,j})=j$ defines a strong $d$-coloring $c\in\Cc_d(H)$, and the distinct vertices in each column $j\in\{0,\ldots,d-1\}$,
\begin{equation}\label{eq:col-state}
S_j=\{M_{i,j}:1\le i\le n\}\subseteq V,
\end{equation}
form $d$ pairwise disjoint two-valued state supports $s_j=\mathbf{1}_{S_j}\in\Sc_{\mathrm{ext}}(H)$ partitioning $V$. Moreover, the vertex degrees and column support sizes obey the exact conservation laws
\begin{align}
\sum_{v\in S_j}\deg_H(v) &= n \qquad (0\le j\le d-1),\label{eq:deg-conservation}\\
\sum_{j=0}^{d-1}|S_j| &= |V|,\label{eq:size-conservation}\\
|S_j| &= n - \sum_{v\in S_j}\bigl(\deg_H(v)-1\bigr).\label{eq:hinge-saving}
\end{align}
\end{theorem}
\begin{proof}
Given $c\in\Cc_d(H)$, each context $C_i$ contains a unique vertex $v\in C_i$ of color $c(v)=j$; setting $M_{i,j}=v$ satisfies Row Bijectivity, and since $c(v)$ is well-defined on $V$, $v$ appears only in column $j=c(v)$, establishing Eq.~\eqref{eq:col-alignment}. Conversely, given $M$ satisfying (1) and (2), Eq.~\eqref{eq:col-alignment} assigns a unique column index $c(v)=j$ to every $v\in V=\bigcup_{i=1}^n C_i$, and Row Bijectivity ensures that $c$ is injective on each $C_i$. Because column $j$ contains one entry from each of the $n$ rows and shares no vertices with any column $j'\neq j$, its underlying vertex set $S_j$ hits every context $C_i$ exactly once, so $\mathbf{1}_{S_j}\in\Sc_{\mathrm{ext}}(H)$. Each vertex $v\in S_j$ appears in $\deg_H(v)$ rows of column $j$, summing to $n$ entries in column $j$ [Eq.~\eqref{eq:deg-conservation}] and yielding Eq.~\eqref{eq:hinge-saving}; summing $|S_j|$ over the $d$ disjoint columns partitioning $V$ gives Eq.~\eqref{eq:size-conservation}.
\end{proof}

\begin{remark}[Incidence equations and wider palettes]\label{rem:wider-palette}
Let $B\in\{0,1\}^{n\times |V|}$ be the incidence matrix and let $X_{v,a}=1$ mean that $v$ has color $a$. The criterion is equivalently
\begin{equation}\label{eq:incidence-factorization}
X\mathbf1_d=\mathbf1_{|V|},\qquad BX=J_{n,d},\qquad X\in\{0,1\}^{|V|\times d}.
\end{equation}
If a palette of size $k>d$ is used, the corresponding $n\times k$ vertex matrix has $k-d$ blanks in each row, and its column supports are independent sets meeting each context \emph{at most} once. They need not be states. Algebraically $X\mathbf1_k=\mathbf1_{|V|}$ and $BX\le J_{n,k}$. The state-column interpretation applies to the exact palette $k=d$.
\end{remark}


The matrix describes the spectral labels on already specified quantum projectors. It does not reconstruct their coordinates. If $u\sim_c v$, meaning $c(u)=c(v)$ in every $c\in\Cc_d(H)$, the corresponding rays remain distinct; it is their allowed scalar labels that coincide. Likewise, a fixed state column may fail to admit the remaining $d-1$ columns, producing a nonextendable state of the realized measurement scenario.

\subsection{A physical two-context schematic}\label{sec:L12}
Take the real rays
\begin{equation}\label{eq:physical-firefly-vectors}
\begin{aligned}
r_1&=(1,0,0),&r_2&=(0,1,0),&r_3&=(0,0,1),\\
r_4&=(1,1,0),&r_5&=(1,-1,0).
\end{aligned}
\end{equation}
Set $q_v=r_v/\|r_v\|$. The contexts $C_1=\{1,2,3\}$ and $C_2=\{3,4,5\}$ are orthonormal bases, sharing projector $P_3$. Direct dot products verify faithfulness. One consistent spectral assignment is
\begin{equation}\label{eq:physical-firefly-matrix}
M=\begin{pmatrix}1&2&3\\4&5&3\end{pmatrix},\qquad
\begin{gathered}
S_0=\{1,4\},\quad S_1=\{2,5\},\\S_2=\{3\}.
\end{gathered}
\end{equation}
Thus both observables use the same spectrum, and the shared projector has eigenvalue $\lambda_2$ in both. Figure~\ref{fig:matrix-reading} shows the two actual contexts; the repeated entry $3$ in Eq.~\eqref{eq:physical-firefly-matrix} represents one projector. This is the minimal measurement example needed to read the larger matrices. In all hypergraph drawings, plane positions are diagram coordinates; Hilbert-space orthogonality is certified by the supplied ray vectors.
\begin{figure}[t]
\centering
\input{2026-latin-context-matrix-fig-firefly.tikz}
\caption{Two projective measurements in $\mathbb R^3$ with the coordinates~\eqref{eq:physical-firefly-vectors}. Each numbered ray occurs once, and each colored line is a complete context. The two incidence rings at vertex $3$ mark the shared projector. The associated $2\times3$ matrix is Eq.~\eqref{eq:physical-firefly-matrix}; line colors distinguish measurements, not eigenvalue labels.}\label{fig:matrix-reading}
\end{figure}

\subsection{Latin indexing and physical color signatures}\label{sec:row-col}\label{sec:graeco-latin}

When a coordinatized scenario is indexed by a $d\times d$ cell array and its contexts include every row and column, a strong $d$-coloring is a Latin square: each spectral label occurs once per row and column. A binary state selects one cell in every row and column and is therefore a permutation transversal. If the actual measurement family also contains the symbol classes of a Latin square $L$, states must meet every such class once, and a strong coloring is an orthogonal mate of $L$. This is the incidence dictionary used for the realized net in Sec.~\ref{sec:latin-for}; its coordinates are independently established there.

For arbitrary coordinatized contexts, two chosen colorings define
\begin{equation}\label{eq:signature-map}
\sigma(v)=\bigl(c_1(v),c_2(v)\bigr).
\end{equation}
On a cell-indexed row--column scenario, this signature is injective precisely when the two Latin colorings are orthogonal, meaning that each ordered pair occurs once. On a general quantum scenario it is simply a table of allowed spectral labels. Several distinct projectors may have one signature; Sec.~\ref{sec:H17-H21} gives a physical example. The $n\times d$ matrix itself requires no literal Latin-square embedding.



\section{A realized Latin net for seven-level measurements}\label{sec:latin-for}

Take the cyclic Latin square $L(i,j)=i+j\pmod7$. We construct forty-nine distinct rank-one projectors whose rows, columns, and symbol classes are twenty-one complete measurements of a seven-level quantum system. Put $\omega=e^{2\pi i/7}$ and define the permutation $f$ by
\begin{equation}\label{eq:for7-permutation}
\begin{array}{c|rrrrrrr}
x&0&1&2&3&4&5&6\\\hline
f(x)&0&2&5&1&6&4&3.
\end{array}
\end{equation}
For each $(i,j)\in\mathbb Z_7^2$, set
\begin{equation}\label{eq:for7-coordinates}
q_{ij}=\frac1{\sqrt7}\bigl(\omega^{\,i f(x)+j x}\bigr)_{x=0}^{6}\in\mathbb C^7.
\end{equation}
This is the Fourier quantum-Latin construction used in Ref.~\cite{ball2026large}; the specific permutation and the full faithfulness check below are explicit finite certificates for the present example. For instance, the exponent vectors of $q_{00},q_{01},q_{10}$ are respectively $(0,0,0,0,0,0,0)$, $(0,1,2,3,4,5,6)$, and $(0,2,5,1,6,4,3)$. Equation~\eqref{eq:for7-coordinates} specifies all forty-nine vectors without numerical approximation.

\begin{proposition}[An exact $49$-ray FOR in $\mathbb C^7$]\label{prop:cyclic7-for}
The vectors~\eqref{eq:for7-coordinates} represent $49$ distinct rays and satisfy
\begin{equation}\label{eq:for7-faithfulness}
q_{ij}\perp q_{k\ell}\quad\Longleftrightarrow\quad
\begin{gathered}
i=k\ \text{or}\ j=\ell\\
\text{or}\ i+j=k+\ell\pmod7,
\end{gathered}
\end{equation}
for distinct cells. Thus they form a FOR of the cyclic $7$-uniform $49$-vertex, $21$-context net.
\end{proposition}
\begin{proof}
The inner product is $7^{-1}\sum_x\omega^{a f(x)+b x}$, where $a=k-i$ and $b=\ell-j$. A sum of seven seventh roots of unity vanishes exactly when its seven exponents occur once each: writing their multiplicities as $h_0,\ldots,h_6$, the irreducibility of $1+z+\cdots+z^6$ makes $\sum h_r\omega^r=0$ equivalent to $h_0=\cdots=h_6=1$.

If $a=0$ and $b\ne0$, the exponents are a permutation. If $a\ne0$, put $m=-b/a$. Directly from Eq.~\eqref{eq:for7-permutation},
\[
\bigl(|\operatorname{im}(f-m\,\mathrm{id})|\bigr)_{m=0}^{6}
=(7,7,4,4,4,4,4).
\]
Therefore the inner product vanishes exactly for $m=0$ or $m=1$, namely $b=0$ or $a+b=0$. Together with $a=0$, these are precisely the three listed context directions. This checks nonedges as well as edges.

All first components are $1/\sqrt7$, so proportionality would mean equality. Equal vectors at different cells would give $a f(x)+b x=0$ for every $x$. If $a\ne0$, this would make $f$ linear, whereas $f(1)=2$ and $f(2)=5\ne4$. If $a=0$, then also $b=0$. Thus all rays are distinct.
\end{proof}


For a density operator $\rho$, the probabilities $p_{ij}=\operatorname{Tr}(\rho P_{ij})$ sum to one on every row, column, and symbol context. These are three families of seven laboratory measurement settings sharing identical rank-one projectors at their intersections. The mate $M(i,j)=i+2j$ assigns a common nondegenerate seven-element spectrum to all twenty-one observables via Eq.~\eqref{eq:context-observable}. This example connects the Latin incidence structure to actual Hilbert-space projectors; its purpose is to establish a realized baseline rather than a spectral obstruction.

The certificate checks all $\binom{49}{2}=1176$ pairs exactly: $441$ are orthogonal and the other $735$ are nonorthogonal. Among the supplied rays, the only complete seven-ray bases are the $21$ listed contexts. A classical mate $M(i,j)=i+2j$ gives $\chi=7$; there are $133$ transversal states, and all extend by the group-table argument. Thus this FOR provides a geometric baseline, \emph{not} an example of a chromatic gap or of nonextendable states. It is a complex coordinatization; no real seven-dimensional realization is asserted.

% BEGIN LATIN HYPERGRAPH DRAWINGS cyclicfor
Figure~\ref{fig:cyclic-for-hypergraph} draws this coordinatized net with the same numbered-node and context-ring convention as the earlier figures.
\begin{figure*}[t]
\centering
\input{2026-latin-context-matrix-fig-cyclicfor.tikz}
\caption{Twenty-one complete projective measurements on $49$ distinct rays in $\mathbb C^7$, with coordinates~\eqref{eq:for7-coordinates}. Vertex $v_{7i+j+1}$ represents $q_{ij}$ and projector $P_{ij}$. Horizontal, vertical, and curved contexts fix $i$, $j$, and $i+j\pmod7$, respectively. Proposition~\ref{prop:cyclic7-for} proves exact faithfulness, including nonorthogonality of every unjoined pair. Each complete seven-ray basis among these rays is listed. The configuration admits a common seven-value spectrum and all $133$ ordinary states extend. Context colors identify measurements, not eigenvalue labels.}\label{fig:cyclic-for-hypergraph}
\end{figure*}



\section{Realized quantum configurations with an exact common spectrum}\label{sec:exact-palette}

The following configurations have faithful real coordinates in dimensions three or four. We use their context matrices to analyze global spectral constraints and the distinction between ordinary states and color-class states. Full coordinate inventories and geometric checks for the companion configurations are given in Ref.~\cite{svozil2026contexts,supplement}.

\subsection{The Pentagon \texorpdfstring{$P_5$}{P5} (\texorpdfstring{$5\times 3$}{5x3} Matrix and Purely Chromatic Ternary Exclusion)}\label{sec:P5}

The pentagon hypergraph $P_5$ ($d=3,n=5$ contexts, $|V|=10$, Berge girth $g=5$, real FOR in $\R^3$) has contexts
\begin{equation}\label{eq:P5-contexts}
\begin{aligned}
C_1&=\{1,6,2\},\quad C_2=\{2,7,3\},\quad C_3=\{3,8,4\},\\
C_4&=\{4,9,5\},\quad C_5=\{5,10,1\},
\end{aligned}
\end{equation}
One explicit FOR uses the following unnormalized real vectors, with $q_v=r_v/\|r_v\|$:
\begin{equation}\label{eq:P5-physical-vectors}
\begin{aligned}
r_1&=(1,0,0),&r_6&=(0,1,-1),\\
r_2&=(0,1,1),&r_7&=(2,-1,1),\\
r_3&=(1,1,-1),&r_8&=(1,0,1),\\
r_4&=(1,-2,-1),&r_9&=(5,2,1),\\
r_5&=(0,1,-2),&r_{10}&=(0,2,1).
\end{aligned}
\end{equation}
Their orthogonal pairs are exactly those of the five listed contexts. Figure~\ref{fig:physical-pentagon} therefore depicts five projective measurements, with each shared ray represented once.
\begin{figure}[t]
\centering
\input{2026-latin-context-matrix-fig-pentagon.tikz}
\caption{The real three-dimensional pentagon measurement scenario, with exact coordinates~\eqref{eq:P5-physical-vectors}. Each colored line is an orthonormal triad. Numbered circles are the ten distinct physical rays; concentric rings show shared outcomes. Black rims identify the terminals $6,7,9$, which cannot all receive one spectral label in any consistent three-value assignment. The rays and contexts remain well defined independently of that spectral restriction.}\label{fig:physical-pentagon}
\end{figure}


There are five degree-two cycle hinges $\{1,2,3,4,5\}$ and five degree-one midpoint vertices $\{6,7,8,9,10\}$.

In a $5\times 3$ Context--Color Matrix $M\in V^{5\times 3}$, each hinge $h\in\{1,2,3,4,5\}$ has $\deg(h)=2$ and saves $1$ vertex when placed in a column [Eq.~\eqref{eq:hinge-saving}]. Because the $5$ hinges form an odd $5$-cycle, at most $\lfloor 5/2\rfloor=2$ nonadjacent hinges can share a single column!
\begin{itemize}
\item \emph{Why $s_{11}=\{6,7,8,9,10\}$ is nonextendable:} The state $s_{11}$ uses $k=0$ hinges (size $|S|=5$). If $s_{11}$ is placed in Column~$\mathbf{c}_0$, the remaining $5$ hinges $\{1,2,3,4,5\}$ must be distributed between only $2$ columns ($\mathbf{c}_1,\mathbf{c}_2$). Since each column can hold at most $2$ hinges, two columns can hold at most $2+2=4<5$ hinges—an immediate Pigeonhole obstruction!
\item \emph{Why all other $10$ states extend ($5$ normalized $5\times3$ matrices):} Every valid $5\times 3$ matrix $M$ must distribute the $5$ hinges across the $3$ columns as $2+2+1=5$ hinges, yielding column support sizes $(5-2)+(5-2)+(5-1)=3+3+4=10=|V|$.
Fixing the anchor context $C_1=(1,6,2)$ in Row~1 as $(\mathbf{c}_0,\mathbf{c}_1,\mathbf{c}_2)=(0,1,2)$ gives the $N_c=5$ normalized $5\times 3$ Context--Color Matrices (corresponding to the $5$ choices of the uniquely colored hinge):
\end{itemize}
\begin{equation}\label{eq:P5-matrices}
\begin{aligned}
M_1^{(P_5)}&=\begin{pmatrix}
\mathbf{1}&6&\mathbf{2}\\
\mathbf{3}&7&\mathbf{2}\\
\mathbf{3}&\mathbf{4}&8\\
9&\mathbf{4}&\mathbf{5}\\
\mathbf{1}&10&\mathbf{5}
\end{pmatrix},&
M_2^{(P_5)}&=\begin{pmatrix}
\mathbf{1}&6&\mathbf{2}\\
\mathbf{3}&7&\mathbf{2}\\
\mathbf{3}&8&\mathbf{4}\\
9&\mathbf{5}&\mathbf{4}\\
\mathbf{1}&\mathbf{5}&10
\end{pmatrix},\\[4pt]
M_3^{(P_5)}&=\begin{pmatrix}
\mathbf{1}&6&\mathbf{2}\\
7&\mathbf{3}&\mathbf{2}\\
\mathbf{4}&\mathbf{3}&8\\
\mathbf{4}&\mathbf{5}&9\\
\mathbf{1}&\mathbf{5}&10
\end{pmatrix},&
M_4^{(P_5)}&=\begin{pmatrix}
\mathbf{1}&6&\mathbf{2}\\
7&\mathbf{3}&\mathbf{2}\\
\mathbf{4}&\mathbf{3}&8\\
\mathbf{4}&9&\mathbf{5}\\
\mathbf{1}&10&\mathbf{5}
\end{pmatrix},\\[4pt]
M_5^{(P_5)}&=\begin{pmatrix}
\mathbf{1}&6&\mathbf{2}\\
7&\mathbf{3}&\mathbf{2}\\
8&\mathbf{3}&\mathbf{4}\\
9&\mathbf{5}&\mathbf{4}\\
\mathbf{1}&\mathbf{5}&10
\end{pmatrix}.
\end{aligned}
\end{equation}
Inspection of Eq.~\eqref{eq:P5-matrices} immediately proves the \emph{purely chromatic ternary exclusion} $\neg(c(6)=c(7)=c(9))$: in any $5\times 3$ matrix where $6,7,9$ all sit in Column~$\mathbf{c}_1$, rows $C_1,C_2,C_4$ force the five hinges $\{1,2,3,4,5\}$ entirely into the two remaining columns $\mathbf{c}_0$ and $\mathbf{c}_2$, which can hold at most $4$ hinges!

\subsection{Faithful Yu--Oh Restrictions \texorpdfstring{$H_{17,10}$ and $H_{21,13}$}{H17,10 and H21,13} and the Graeco-Latin Signature Table}\label{sec:H17-H21}

The $17$-ray, $10$-context hypergraph $H_{17,10}$ retains the four base contexts $C_1=\{1,2,3\}$, $C_2=\{7,8,9\}$, $C_3=\{4,5,6\}$, $C_4=\{2,8,5\}$ of completed Yu--Oh $H_{\YO}^+$ together with the six contexts of diagonal sectors $h_1=11$ and $h_2=10$:
\begin{equation}\label{eq:H17-sectors}
\begin{aligned}
\mathcal{E}_{11}&=\bigl\{\{3,11,14\},\{7,11,15\},\{4,11,24\}\bigr\},\\
\mathcal{E}_{10}&=\bigl\{\{9,10,23\},\{1,10,19\},\{4,10,18\}\bigr\}.
\end{aligned}
\end{equation}
Fixing the anchor row $C_4=(2,8,5)$ across columns $(\mathbf{c}_0,\mathbf{c}_1,\mathbf{c}_2)$, $H_{17,10}$ admits $N_c=4$ normalized $10\times 3$ Context--Color Matrices (Table~S41 of Ref.~\cite{supplement}), shown here for $c_1$ and $c_2$:
\begin{equation}\label{eq:H17-matrices}
\setlength{\arraycolsep}{2.2pt}
M_1^{(H_{17,10})}=\begin{pmatrix}
\mathbf{c}_0&\mathbf{c}_1&\mathbf{c}_2\\\hline
2&1&3\\
9&8&7\\
4&6&5\\
2&8&5\\\hline
14&\mathbf{11}&3\\
15&\mathbf{11}&7\\
4&\mathbf{11}&\mathbf{24}\\\hline
9&23&\mathbf{10}\\
19&1&\mathbf{10}\\
4&\mathbf{18}&\mathbf{10}
\end{pmatrix},\
M_2^{(H_{17,10})}=\begin{pmatrix}
\mathbf{c}_0&\mathbf{c}_1&\mathbf{c}_2\\\hline
2&1&3\\
9&8&7\\
6&4&5\\
2&8&5\\\hline
\mathbf{11}&14&3\\
\mathbf{11}&15&7\\
\mathbf{11}&4&\mathbf{24}\\\hline
9&23&\mathbf{10}\\
19&1&\mathbf{10}\\
\mathbf{18}&4&\mathbf{10}
\end{pmatrix}.
\end{equation}
In every coloring of $H_{17,10}$, the Specker-bug restriction forces $c(10)\ne c(11)$ (Sec.~\ref{sec:3d-gadgets}). Both hubs are also adjacent to vertex $4$, through the roof contexts $\{4,11,24\}$ and $\{4,10,18\}$. Thus $4,10,11$ always have three distinct colors. The first roof forces $24$ into the color of $10$, and the second forces $18$ into the color of $11$. Hence $10\sim_c24$ and $11\sim_c18$ in all four normalized matrices.


\paragraph{A direct quantum-probability consequence.}
The real coordinates of these two distinct rays are $r_{10}=(1,-1,1)$ and $r_{24}=(1,-1,2)$. Prepare $\rho=P_{24}=r_{24}r_{24}^{T}/6$. The Born rule gives
\begin{equation}\label{eq:H17-physical-born}
p_{24}=1,\qquad p_{10}=\frac{|r_{10}\cdot r_{24}|^2}{\|r_{10}\|^2\|r_{24}\|^2}=\frac89.
\end{equation}
Every color-class state has $s(10)=s(24)$, so every convex mixture of extendable states has $p_{10}=p_{24}$. This physically available preparation lies outside that convex hull. In any decomposition into ordinary context states, a nonextendable total weight $\eta$ must satisfy $|p_{24}-p_{10}|\le\eta$, and hence $\eta\ge1/9$. The companion supplies the complete Born vector and an ordinary-state mixture attaining this bound~\cite{svozil2026contexts,supplement}. A forced scalar identity therefore does not identify the two quantum projectors.

Adjoining the third diagonal sector $\mathcal{E}_{13}=\{\{1,13,20\},\{7,13,21\},\{6,13,22\}\}$ yields $H_{21,13}$ ($n=13$ contexts, $|V|=21$). Only two normalized $13\times 3$ Context--Color Matrices survive ($M_1^{(H_{21,13})}$ and $M_2^{(H_{21,13})}$, corresponding to the antipodal base patterns $101$ and $010$):
\begin{equation}\label{eq:H21-matrices}
\setlength{\arraycolsep}{2.2pt}
M_1^{(H_{21,13})}=\begin{pmatrix}
\mathbf{c}_0&\mathbf{c}_1&\mathbf{c}_2\\\hline
2&1&3\\
9&8&7\\
4&6&5\\
2&8&5\\\hline
14&11&3\\
15&11&7\\
4&11&24\\\hline
9&23&10\\
19&1&10\\
4&18&10\\\hline
13&1&20\\
13&21&7\\
13&6&22
\end{pmatrix},\
M_2^{(H_{21,13})}=\begin{pmatrix}
\mathbf{c}_0&\mathbf{c}_1&\mathbf{c}_2\\\hline
2&3&1\\
7&8&9\\
6&4&5\\
2&8&5\\\hline
14&3&11\\
7&15&11\\
24&4&11\\\hline
10&23&9\\
10&19&1\\
10&4&18\\\hline
20&13&1\\
7&13&21\\
6&13&22
\end{pmatrix}.
\end{equation}

The joint-signature criterion of Sec.~\ref{sec:graeco-latin} applies as follows: each vertex $v\in V(H_{21,13})$ receives a $2$-tuple chromatic signature $\sigma(v)=(c_1(v),c_2(v))\in\{0,1,2\}^2$ recording its column index in $M_1^{(H_{21,13})}$ and $M_2^{(H_{21,13})}$. Placing the $21$ vertices into the $3\times 3$ signature grid $\{0,1,2\}\times\{0,1,2\}$ yields the $9$ chromatic equivalence classes:
\begin{equation}\label{eq:H21-signature-grid}
\begin{array}{c|ccc}
c_2\backslash c_1 & \mathbf{c}_0=0 & \mathbf{c}_1=1 & \mathbf{c}_2=2\\\hline
0 & \{2,14\} & \{6\} & \{7,10,20,24\}\\
1 & \{4,13,15,19\} & \{8,23\} & \{3\}\\
2 & \{9\} & \{1,11,18,21\} & \{5,22\}
\end{array}
\end{equation}
Whereas a $3\times 3$ Graeco-Latin square has $9$ vertices and places \emph{one} vertex in each of the $9$ cells of Eq.~\eqref{eq:H21-signature-grid}, $H_{21,13}$ has $21$ vertices and only two normalized colorings, forcing multiple vertices to collide in six of the nine cells and producing $3\binom{4}{2}+3\binom{2}{2}=21$ chromatically inseparable pairs. Each pair consists of distinct faithfully represented quantum rays; the common signature constrains their scalar labels, not their coordinates.

\subsection{Three-Dimensional Logical Gadgets \texorpdfstring{$B_{13,7},D_{16,9},I_{16,9},T_{19,11}$}{B13,7, D16,9, I16,9, T19,11}}\label{sec:3d-gadgets}

The remaining three-dimensional gadgets in the companion inventory~\cite{svozil2026contexts} are sub-hypergraphs or roof extensions of $H_{17,10}$, and their $n\times 3$ Context--Color Matrices directly exhibit how TIFS and TITS constraints operate:

\begin{enumerate}
\item \emph{Completed Specker Bug $B_{13,7}$ ($7\times 3$ matrix, $N_c=6$):} Deleting $C_3=\{4,5,6\}$, $\{4,11,24\}$, and $\{4,10,18\}$ from $H_{17,10}$ leaves the $7$ contexts of $B_{13,7}$:
\[
\begin{aligned}
&\{1,2,3\},\ \{7,8,9\},\ \{2,8,5\},\ \{3,11,14\},\\
&\{7,11,15\},\ \{9,10,23\},\ \{1,10,19\}.
\end{aligned}
\]
Suppose $10$ and $11$ had one common color $h$. Their four neighbors $1,3,7,9$ would all avoid $h$. If $h\ne c(2)$, the context $\{1,2,3\}$ could not contain $h$, contrary to its being rainbow. If $h=c(2)$, then $h\ne c(8)$ because $2,8$ share a context; now $\{7,8,9\}$ cannot contain $h$. Both cases contradict strong coloring. Thus $c(10)\ne c(11)$, without making a palette assumption incompatible with the fixed anchor.
\item \emph{Completed $\Gamma_1$ Gadget $I_{16,9}$ ($9\times 3$ matrix, $N_c=6$):} Re-adjoining the two roof contexts $\{4,11,24\}$ and $\{4,10,18\}$ to $B_{13,7}$ (without $\{4,5,6\}$) appends two rows to each of the $6$ matrices of $B_{13,7}$. Since $10$ and $11$ always occupy distinct columns in $B_{13,7}$, the apex $4$ is forced into the third column, locking $24$ into the column of $10$ ($10\sim_c 24$) and $18$ into the column of $11$ ($11\sim_c 18$).
\item \emph{Purely Chromatic Binary TIFS $D_{16,9}$ ($9\times 3$ matrix, $N_c=7$) and Two-Way TITS $T_{19,11}$ ($11\times 3$ matrix, $N_c=7$):} Deleting $\{1,10,19\}$ from $H_{17,10}$ yields $D_{16,9}$. While the nonextendable state $s=\{1,5,7,14,18,23,24\}\in\Sc(D_{16,9})$ contains both $14$ and $18$, no $9\times 3$ Context--Color Matrix can place $14$ and $18$ in the same column: if $14,18\in\mathbf{c}_0$, then $\{3,11,14\},\{4,11,24\},\{4,10,18\}$ force $3,4$ into one column (say $\mathbf{c}_1$) and $10,11$ into $\mathbf{c}_2$, which forces $7,9\notin\mathbf{c}_2\implies 8\in\mathbf{c}_2$, and $2,5\notin\mathbf{c}_1\implies 2,5\in\mathbf{c}_0$, colliding in Row~4 ($\{2,8,5\}$)! Adjoining the two roof contexts $\{14,26,27\}$ and $\{18,26,28\}$ in $T_{19,11}$ appends two rows to the $7$ normalized $9\times 3$ matrices of $D_{16,9}$; since $14$ and $18$ always sit in distinct columns, their common neighbor $26$ must occupy the third column, locking $28$ into the column of $14$ ($14\sim_c 28$) and $27$ into the column of $18$ ($18\sim_c 27$).
\end{enumerate}

\subsection{Four-Dimensional Apex Lifts and Strictly Linear Gadgets (\texorpdfstring{$n\times 4$}{n x 4} Matrices)}\label{sec:4d-gadgets}

In dimension $d=4$, the $n\times 4$ Context--Color Matrix separates into two distinct structural classes (the companion four-dimensional inventory~\cite{svozil2026contexts}):

\begin{proposition}[Common-apex lift in arbitrary dimension]\label{prop:general-apex}
Let $H$ be $d_0$-uniform and adjoin $r=d-d_0$ common fresh vertices $A$ to every context. Then
\[
\Gamma(H^A)=\Gamma(H)\vee K_r,\qquad\chi(H^A)=\chi(H)+r.
\]
Its states are the old states extended by zero on $A$, together with the $r$ singleton apex states. If $\chi(H)=d_0$, its extendable states transform in the same way, so the old nonextendable states and forced color identities persist. Any injective faithful realization in $\F^{d_0}$ extends faithfully into $\F^d$ by adjoining coordinate rays.
\end{proposition}
\begin{proof}
All new vertices are adjacent to each other and to every old vertex, proving the graph join and chromatic formula. A state either selects one apex, forcing all old vertices to zero, or selects no apex and restricts to an old state. In an exact-palette coloring the apices use $r$ distinct colors absent from the old vertices, leaving precisely $d_0$ colors on $H$. Explicitly, embed each old vector as $(q_v,0_r)$ and use $e_{d_0+1},\ldots,e_d$ for the apices. Every enlarged context is an orthonormal basis, and the construction preserves exactly the old orthogonalities and adds the required apex pairs.
\end{proof}


\subsubsection{Non-Linear Common-Apex Lifts \texorpdfstring{$W_{11,5}^{(4)}$, $B_{14,7}^{(4)}$, and $D_{17,9}^{(4)}$}{W11,5(4), B14,7(4), and D17,9(4)}}
Each common-apex lift adjoins a single apex vertex $0$ to every $3$-element context $C_i$ of a $3$-dimensional gadget $H$, forming $C_i^{(4)}=\{0\}\cup C_i$. In Theorem~\ref{thm:nxd-matrix}, since $0\in\bigcap_{i=1}^n C_i^{(4)}$ has degree $\deg(0)=n$, Eq.~\eqref{eq:hinge-saving} gives $|S_0|=n-(n-1)=1$. Thus \emph{every} $n\times 4$ Context--Color Matrix $M^{(4)}$ of an apex lift is obtained by prepending a constant Column~$\mathbf{c}_0=(0,0,\ldots,0)^T$ to a valid $n\times 3$ Context--Color Matrix $M^{(3)}$ of $H$:
\begin{equation}\label{eq:apex-matrix}
M^{(4)}=\left(\begin{array}{c|ccc}
\mathbf{c}_0 & \mathbf{c}_1 & \mathbf{c}_2 & \mathbf{c}_3\\\hline
0 & & & \\
0 & & M^{(3)} & \\
\vdots & & & \\
0 & & &
\end{array}\right)\in V^{n\times 4}.
\end{equation}
Equation~\eqref{eq:apex-matrix} provides a one-line matrix proof of the apex-extension proposition of Ref.~\cite{svozil2026contexts}: the normalized coloring count $N_c$ is unchanged, the singleton column $S_0=\{0\}$ is always an extendable state, and the remaining three columns $\mathbf{c}_1,\mathbf{c}_2,\mathbf{c}_3$ are identical to those of $M^{(3)}$.

\subsubsection{Strictly Linear Four-Dimensional Gadgets \texorpdfstring{$W_{16,6}^{(4)}$, $L_{18,7}^{(4)}$, and $L_{22,9}^{(4)}$}{W16,6(4), L18,7(4), and L22,9(4)}}
To restore linearity ($|e\cap e'|\le 1$, so $g(H)\ge 3$), Ref.~\cite{svozil2026contexts} constructs three faithful linear $4$-uniform gadgets in $\R^4$:
\begin{itemize}
\item \emph{Linear $5$-Wheel $W_{16,6}^{(4)}$ ($6\times 4$ matrix, $N_c=80$):} The $6$ contexts are
\[
\begin{aligned}
C_1&=\{2,0,6,1\},\quad C_2=\{2,7,13,3\},\\
C_3&=\{3,0,8,4\},\quad C_4=\{4,9,14,5\},\\
C_5&=\{5,10,15,1\},\quad C_6=\{0,11,12,5\}.
\end{aligned}
\]
Here vertex $0$ lies in only three contexts ($C_1,C_3,C_6$). In any $6\times 4$ matrix $M$, placing $0$ in Column~$\mathbf{c}_3$ covers $0$ in Rows~$1,3,6$. If the three terminals $6,8,11$ were all placed in Column~$\mathbf{c}_0$, then in Rows~$1,3,6$ the outer cycle vertices $\{1,2,3,4,5\}$ would be barred from both $\mathbf{c}_0$ and $\mathbf{c}_3$! Thus $\{1,2,3,4,5\}$ would have to be distributed across only two columns ($\mathbf{c}_1,\mathbf{c}_2$) while pairwise sharing rows around the $5$-cycle $C_1-C_2-C_3-C_4-C_5$, which is impossible. One explicit $6\times4$ matrix is:
\begin{equation}\label{eq:W16-matrix}
M_1^{(W_{16,6}^{(4)})}=\begin{pmatrix}
\mathbf{c}_0 & \mathbf{c}_1 & \mathbf{c}_2 & \mathbf{c}_3\\\hline
2 & \mathbf{0} & \mathbf{6} & 1\\
2 & 7 & 3 & 13\\
\mathbf{8} & \mathbf{0} & 3 & 4\\
5 & 9 & 14 & 4\\
5 & 10 & 15 & 1\\
5 & \mathbf{0} & \mathbf{11} & 12
\end{pmatrix}.
\end{equation}
\item \emph{Linear Quaternion TIFS $L_{18,7}^{(4)}$ ($7\times 4$ matrix, $N_c=144$, $\Sc_{\mathrm{ext}}=\Sc=68$) and Linear TITS $L_{22,9}^{(4)}$ ($9\times 4$ matrix, $N_c=216$):} In $L_{18,7}^{(4)}$, the $7$ contexts are
\[
\begin{aligned}
C_1&=\{1,2,4,3\},\ C_2=\{5,8,7,6\},\ C_3=\{4,9,10,8\},\\
C_4&=\{11,1,12,5\},\ C_5=\{11,14,13,2\},\\
C_6&=\{15,3,16,7\},\ C_7=\{15,17,18,6\}.
\end{aligned}
\]
Why can $11$ and $15$ never share a column in a $7\times 4$ matrix (or even a single state column)? If $11,15\in\mathbf{c}_0$, then Rows~$4,5,6,7$ bar $\{1,2,3\}$ and $\{5,6,7\}$ from $\mathbf{c}_0$; then Row~1 ($\{1,2,4,3\}$) forces $4\in\mathbf{c}_0$ and Row~2 ($\{5,8,7,6\}$) forces $8\in\mathbf{c}_0$, so $4$ and $8$ collide in Column~$\mathbf{c}_0$ of Row~3 ($\{4,9,10,8\}$)! In $L_{22,9}^{(4)}$, adjoining $C_8=\{11,21,22,19\}$ and $C_9=\{13,20,15,19\}$ whose listed context $C_9$ includes the pair $\{13,15\}$ ensures that $\{11,13,15,19\}$ occupy all four distinct columns, locking $20\in C_9$ into the column of $11$ ($11\sim_c 20$). An explicit $9\times 4$ Context--Color Matrix (with its first context fixed as $(1,2,3,4)$) is:
\begin{equation}\label{eq:L22-matrix}
M_1^{(L_{22,9}^{(4)})}=\begin{pmatrix}
\mathbf{c}_0 & \mathbf{c}_1 & \mathbf{c}_2 & \mathbf{c}_3\\\hline
1 & 2 & 3 & 4\\
7 & 5 & 8 & 6\\
9 & 10 & 8 & 4\\
1 & 5 & \mathbf{11} & 12\\
13 & 2 & \mathbf{11} & 14\\
7 & 15 & 3 & 16\\
17 & 15 & 18 & 6\\
21 & 22 & \mathbf{11} & 19\\
13 & 15 & \mathbf{20} & 19
\end{pmatrix}.
\end{equation}
Notice how Column~$\mathbf{c}_2$ in Eq.~\eqref{eq:L22-matrix} locks $\mathbf{11}$ (in Rows~$4,5,8$) and $\mathbf{20}$ (in Row~$9$) into the same state column $S_2=\{3,8,11,18,20\}$.
\end{itemize}

% ============================================================================


\section{Realized quantum spectral obstructions}\label{sec:obstructions}

The configurations in this section have actual faithful quantum coordinates even though no exact-palette context--color matrix exists. Their spectral inconsistency is compatible with the quantum measurements themselves. Coordinates and all context inventories are supplied in Ref.~\cite{svozil2026contexts,supplement}.

\subsection{Completed Yu--Oh \texorpdfstring{$H_{\YO}^+$}{HYO+} (\texorpdfstring{$16\times 3$}{16x3} Obstruction) and \texorpdfstring{$H_{\YO}^{(4)}$}{HYO(4)}}\label{sec:HYO}

Completed Yu--Oh $H_{\YO}^+$ has $n=16$ contexts of size $d=3$ on $|V|=25$ vertices: the $4\times 3$ base block $E_{\mathrm{base}}=\{\{1,2,3\},\{7,8,9\},\{4,5,6\},\{2,8,5\}\}$ and four $3\times 3$ diagonal sector blocks $\mathcal{E}_{h}$ ($h\in\{10,11,12,13\}$), where each $\mathcal{E}_h$ consists of three rows $\{x,h,w_{xh}\}$ for $x\in N(h)$.

In a $16\times 3$ Context--Color Matrix $M$:
\begin{enumerate}
\item In each $3\times 3$ sector block $\mathcal{E}_h$, the hub $h$ appears in all three rows and must therefore occupy a \emph{single column} $j_h\in\{\mathbf{c}_0,\mathbf{c}_1,\mathbf{c}_2\}$.
\item By Row Bijectivity, none of the three neighbors $x\in N(h)$ can occupy column $j_h$! That is, the three neighbors $N(h)\subset\{1,3,4,6,7,9\}$ are restricted to the \emph{two remaining columns} $\{0,1,2\}\setminus\{j_h\}$.
\item However, in the $4\times 3$ base block (fixing Row~4 as $(2,8,5)$), there are only $2^3=8$ ways to order the pairs $\{1,3\},\{7,9\},\{4,6\}$ across the two available columns in Rows~1, 2, and 3:
\begin{equation}\label{eq:YO-base-matrix}
M_{\mathrm{base}}(b_1,b_2,b_3)=\begin{pmatrix}
\mathbf{c}_0 & \mathbf{c}_1 & \mathbf{c}_2\\\hline
2 & 1\text{ or }3 & 3\text{ or }1\\
9\text{ or }7 & 8 & 7\text{ or }9\\
6\text{ or }4 & 4\text{ or }6 & 5\\
2 & 8 & 5
\end{pmatrix}.
\end{equation}
As shown in the base-coloring partition theorem of Ref.~\cite{svozil2026contexts}, each of the $8$ choices of $(b_1,b_2,b_3)\in\{0,1\}^3$ places one of the four triples $N(10)=\{1,9,4\}$, $N(11)=\{3,7,4\}$, $N(12)=\{3,9,6\}$, or $N(13)=\{1,7,6\}$ across \emph{all three columns} $\mathbf{c}_0,\mathbf{c}_1,\mathbf{c}_2$—leaving \emph{zero} available columns for the corresponding hub $h_j$!
\end{enumerate}
Consequently, no $16\times 3$ Context--Color Matrix exists ($\chi(H_{\YO}^+)=4$). Appending the common apex $26$ in $H_{\YO}^{(4)}$ prepends the singleton column $\mathbf{c}_0=(26,\ldots,26)^T$ as in Eq.~\eqref{eq:apex-matrix}, shifting the exact same obstruction to $16\times 4$ ($\chi(H_{\YO}^{(4)})=5$).

\subsection{The \texorpdfstring{$37$}{37}-Ray Complex Hesse Hypergraph \texorpdfstring{$H_{37}$}{H37}: An Affine-Plane / MOLS Obstruction}\label{sec:H37}

The $37$-ray, $24$-context complex Hesse hypergraph $H_{37}$ in $\C^3$ (the Hesse section of Ref.~\cite{svozil2026contexts}) provides a context hypergraph organized by a finite-geometry Latin-square / MOLS skeleton.

Consider the $9$ core vertices $M_{ak}$ ($a,k\in\Z_3$), arranged as a $3\times 3$ grid:
\begin{equation}\label{eq:H37-grid}
\begin{array}{c|ccc}
M_{ak} & k=0 & k=1 & k=2\\\hline
a=0\ (\mathcal{M}_0) & M_{00}\ (2) & M_{01}\ (3) & M_{02}\ (4)\\
a=1\ (\mathcal{M}_1) & M_{10}\ (5) & M_{11}\ (6) & M_{12}\ (7)\\
a=2\ (\mathcal{M}_2) & M_{20}\ (8) & M_{21}\ (9) & M_{22}\ (10)
\end{array}
\end{equation}
The grid columns $N(P_k)$ and cyclic symbol classes $N(Q_k)$ below are neighborhoods, not contexts. They must omit a color, whereas the actual row contexts must be rainbow. Thus these patterns do not make the core into a $3$-net hypergraph $H_L$.

Look at how the $24$ contexts of $H_{37}$ [the listed Hesse contexts in Ref.~\cite{svozil2026contexts}] act on this $3\times 3$ grid in a putative $24\times 3$ Context--Color Matrix:
\begin{enumerate}
\item \emph{Rows of the $3\times 3$ Grid ($\mathcal{M}_a=\{M_{a0},M_{a1},M_{a2}\}$):} The first $3$ contexts $\mathcal{M}_0,\mathcal{M}_1,\mathcal{M}_2$ are the \emph{rows} of Eq.~\eqref{eq:H37-grid}, forcing each row of Eq.~\eqref{eq:H37-grid} to contain all three colors $\{\mathbf{c}_0,\mathbf{c}_1,\mathbf{c}_2\}$.
\item \emph{Columns of the $3\times 3$ Grid ($\mathcal{B}_{ak}=\{M_{ak},P_k,B_{ak}\}$):} For each fixed $k\in\Z_3$, vertex $P_k$ appears in $\mathcal{B}_{0k},\mathcal{B}_{1k},\mathcal{B}_{2k}$ alongside the $k$th \emph{column} $N(P_k)=\{M_{0k},M_{1k},M_{2k}\}$ of Eq.~\eqref{eq:H37-grid}. Thus column $k$ of the grid~\eqref{eq:H37-grid} must omit the color $c(P_k)$!
\item \emph{Cyclic Latin-Square Transversals ($\mathcal{C}_{ak}=\{M_{a,a+k},Q_k,C_{ak}\}$):} For each fixed $k\in\Z_3$, vertex $Q_k$ appears alongside the diagonal transversal
\[
N(Q_k)=\{M_{0,k},\,M_{1,k+1},\,M_{2,k+2}\},
\]
which is the $k$th symbol class of the cyclic Latin square $L(a,\ell)=\ell-a\pmod 3$! Thus each cyclic diagonal $N(Q_k)$ must omit the color $c(Q_k)$.
\item \emph{The Apex Constraint ($\mathcal{A}_k=\{z,P_k,A_k\}$):} Fix $z$ in Column~$\mathbf{c}_0$ ($c(z)=0$). Then the three rows $\mathcal{A}_0,\mathcal{A}_1,\mathcal{A}_2$ bar $P_0,P_1,P_2$ from Column~$\mathbf{c}_0$, so $c(P_k)\in\{1,2\}$.
\end{enumerate}

Now watch how the Latin-square structure produces the contradiction $\chi(H_{37})=4$:
\begin{itemize}
\item If $P_0,P_1,P_2$ all had the same color (say $1$), then by (2) no entry $M_{ak}$ in any column of Eq.~\eqref{eq:H37-grid} could have color $1$, contradicting (1) that every row $\mathcal{M}_a$ contains all three colors!
\item Hence $\{c(P_0),c(P_1),c(P_2)\}=\{1,2\}$. Up to a cyclic shift of the column index $k\in\Z_3$ and exchanging colors $1,2$, set $(c(P_0),c(P_1),c(P_2))=(2,1,1)$.
\item Since Column~$0$ of Eq.~\eqref{eq:H37-grid} omits color $2$, and Columns~$1$ and~$2$ omit color $1$, the unique color-$1$ entry in each row $\mathcal{M}_a$ \emph{must sit in Column~$0$} ($c(M_{a0})=1$ for all $a\in\Z_3$), while the remaining two entries $(M_{a1},M_{a2})$ in each row must be a permutation $(x_a,2-x_a)$ of $\{0,2\}$:
\begin{equation}\label{eq:H37-forced-grid}
c(M)=\begin{pmatrix}
1 & x_0 & 2-x_0\\
1 & x_1 & 2-x_1\\
1 & x_2 & 2-x_2
\end{pmatrix},\qquad x_a\in\{0,2\}.
\end{equation}
\item Now read off the three cyclic Latin-square transversals $N(Q_k)$ from Eq.~\eqref{eq:H37-forced-grid}:
\[
\begin{aligned}
c(N(Q_0))&=\{1,\,x_1,\,2-x_2\},\\
c(N(Q_1))&=\{x_0,\,2-x_1,\,1\},\\
c(N(Q_2))&=\{2-x_0,\,1,\,x_2\}.
\end{aligned}
\]
By (3), none of these three transversals can be rainbow ($\{0,1,2\}$), so the two non-$1$ entries in each transversal must be equal:
\[
x_1=2-x_2,\qquad x_0=2-x_1,\qquad x_2=2-x_0,
\]
which requires $x_0=2-x_1=2-(2-x_2)=x_2=2-x_0\implies x_0=1\notin\{0,2\}$, a contradiction!
\end{itemize}
Thus no $24\times3$ Context--Color Matrix exists, proving $\chi(H_{37})\ge4$. A strong four-coloring is explicit: set $c(P_k)=c(Q_k)=3$, $c(z)=0$, $c(A_k)=1$, $c(M_{ak})=k$, $c(B_{ak})=k+1\pmod3$, and $c(C_{ak})=a+k+1\pmod3$. Every listed context has three distinct colors. This supplies the upper bound and proves $\chi(H_{37})=4$; its four-column matrix has one blank per row.

\subsection{The Four-Dimensional Completion \texorpdfstring{$H_{65}$}{H65} (\texorpdfstring{$29\times 4$}{29x4} Pigeonhole Overflow)}\label{sec:H65}

Finally, consider the canonical linear $4$-uniform hypergraph $H_{65}$ in $\R^4$ ($|V|=65,n=29$ contexts, the real-completion section of Ref.~\cite{svozil2026contexts}), built from the $17$-ray core $V_{\mathrm{core}}=\{1,\ldots,17\}$ with $5$ core contexts
\begin{equation}\label{eq:H65-core}
\begin{aligned}
C_1&=\{1,2,3,4\},\ C_2=\{5,6,7,8\},\ C_3=\{9,10,12,13\},\\
C_4&=\{9,11,14,15\},\ C_5=\{10,11,16,17\},
\end{aligned}
\end{equation}
and $24$ completed bridge contexts $E_m=\{u_m,v_m,2m+16,2m+17\}$ ($1\le m\le 24$).

In the $29\times 4$ Context--Color Matrix picture of Theorem~\ref{thm:nxd-matrix}, the non-4-colorability $\chi(H_{65})=5$ (despite $H_{65}$ having $|\Sc(H_{65})|=393{,}216$ two-valued states!) has a one-line Pigeonhole proof:
\begin{enumerate}
\item Every state column $S_j$ in a $29\times 4$ Context--Color Matrix must be an independent set of $\Gamma(H_{65})$.
\item On the $17$-vertex core $V_{\mathrm{core}}$, the independence number is $\alpha(G_{\mathrm{core}})=4$: every state $s\in\Sc(H_{65})$ selects \emph{at most $4$ core vertices} (in fact $|S_j\cap V_{\mathrm{core}}|=4$: one from $C_1$, one from $C_2$, one hub from $\{9,10,11\}$, and one bridge ray from $B_1\cup B_2\cup B_3$).
\item Therefore, four disjoint state columns $S_0,S_1,S_2,S_3$ can cover at most
\begin{equation}\label{eq:H65-pigeonhole}
\sum_{j=0}^{3}\bigl|S_j\cap V_{\mathrm{core}}\bigr|\le 4\times 4=16 < 17=\bigl|V_{\mathrm{core}}\bigr|
\end{equation}
core vertices. At least one of the $17$ core vertices is left without a column, proving $\chi(H_{65})\ge5$. The companion supplies a five-coloring, giving $\chi(H_{65})=5$. Such a five-column representation has one blank per row; its columns are independent sets, as in Remark~\ref{rem:wider-palette}.
\end{enumerate}




\section{Physical interpretation and summary}\label{sec:conclusion}

The context--color matrix organizes spectral labels on existing quantum projectors. A valid column partition gives a common nondegenerate spectrum across intertwined measurements; a column obstruction excludes that labeling. The three- and four-dimensional examples retain their quantum coordinates in either case. Table~\ref{tab:master-matrix-synthesis} summarizes the physically realized configurations used here.

\begin{table*}[t]
\centering\small
\caption{Coordinatized quantum configurations discussed in the text. $N_c$ counts exact-palette colorings modulo global palette permutations; a dash means that the total was not enumerated here. $S_{\rm ext}/S$ lists extendable and ordinary state counts. The field and dimension specify the ray realization. Complete companion coordinates are in Ref.~\cite{svozil2026contexts,supplement}; the two-context, pentagon, and cyclic-net coordinates are printed in this paper.}\label{tab:master-matrix-synthesis}
\begin{tabular}{lcrrrr}\toprule
Configuration & Ray space & $|V|$ & $|E|$ & $N_c$ & $|\Sc_{\rm ext}|/|\Sc|$\\\midrule
Two contexts & $\mathbb R^3$ &5&2&2&5/5\\
$P_5$ & $\mathbb R^3$ &10&5&5&10/11\\
$H_{17,10}$ & $\mathbb R^3$ &17&10&4&10/18\\
$H_{21,13}$ & $\mathbb R^3$ &21&13&2&6/21\\
$B_{13,7}$ & $\mathbb R^3$ &13&7&6&12/14\\
$I_{16,9}$ & $\mathbb R^3$ &16&9&6&12/22\\
$D_{16,9}$ & $\mathbb R^3$ &16&9&7&14/21\\
$T_{19,11}$ & $\mathbb R^3$ &19&11&7&14/25\\
$W_{11,5}^{(4)}$ & $\mathbb R^4$ &11&5&5&11/12\\
$B_{14,7}^{(4)}$ & $\mathbb R^4$ &14&7&6&13/15\\
$D_{17,9}^{(4)}$ & $\mathbb R^4$ &17&9&7&15/22\\
$W_{16,6}^{(4)}$ & $\mathbb R^4$ &16&6&80&56/72\\
$L_{18,7}^{(4)}$ & $\mathbb R^4$ &18&7&144&68/68\\
$L_{22,9}^{(4)}$ & $\mathbb R^4$ &22&9&216&92/158\\
Completed Yu--Oh & $\mathbb R^3$ &25&16&0&0/24\\
$H_{\YO}^{(4)}$ & $\mathbb R^4$ &26&16&0&0/25\\
$H_{37}$ & $\mathbb C^3$ &37&24&0&0/153\\
$H_{65}$ & $\mathbb R^4$ &65&29&0&0/393,216\\
Cyclic Latin net & $\mathbb C^7$ &49&21&--&133/133\\\bottomrule
\end{tabular}
\end{table*}

The physical distinction is particularly sharp in $H_{17,10}$: all spectral assignments identify labels $10$ and $24$, while the Born preparation~\eqref{eq:H17-physical-born} distinguishes their probabilities and requires nonextendable ordinary-state weight. The realized Yu--Oh, Hesse, and $65$-ray configurations instead obstruct the exact common spectrum altogether. These are statements about specified quantum measurement families, not deductions of geometry from coloring.

Latin-square structure has a limited, concrete role. It supplies the incidence indexing and common-spectrum coloring of the coordinatized seven-level net, and it clarifies the spectral signatures and neighborhood constraints in the realized three-dimensional examples. The general matrix criterion applies in every Hilbert-space dimension once the quantum ray representation has been supplied. Common-apex extensions preserve that physical foundation through their explicit coordinate embedding.


\begin{acknowledgments}
This research was funded in part by the Austrian Science Fund (FWF), Grant DOI \href{https://doi.org/10.55776/PIN5424624}{10.55776/PIN5424624}.
\end{acknowledgments}


\clearpage
\appendix
\section{Verification of the physical examples}\label{app:verification}

The accompanying certificate contains only the coordinatized examples retained in this paper. Integer dot products verify the full faithful orthogonality relation for Eqs.~\eqref{eq:physical-firefly-vectors} and~\eqref{eq:P5-physical-vectors}, including all nonedges and distinctness of rays. For the cyclic FOR, the checker uses exponent multiplicities modulo seven to verify all $1176$ inner products exactly, with $441$ zeros and $735$ nonzeros, and checks that the only complete seven-ray bases are the $21$ listed measurements. It also enumerates the $133$ transversals. The phase-exponent coordinates are separately provided in \texttt{cyclic7-faithful-coordinates.json}.

For the finite context inventories, normalized coloring fixes one context's color order, giving one representative of each global palette-permutation orbit. Binary states are enumerated with exact-one propagation on the listed measurements. Color-class supports give $\Sc_{\rm ext}$. Each displayed vertex matrix is checked against its context rows and for consistent column positions of shared projectors. Ray coordinates and geometric verification for the companion configurations are those of Ref.~\cite{svozil2026contexts,supplement}; the standalone checker does not claim to rederive those larger coordinate inventories.

The three TikZ files show only coordinatized quantum scenarios. Their trace memberships, unique vertex labels, and incidence rings are checked against the physical context inventories. Trace colors distinguish measurements, while the matrix columns distinguish spectral labels. The standalone \texttt{verify-latin-context-matrix.py} needs only Python's standard library. All active source files, coordinate certificates, and verification output are included in the accompanying package.

\bibliography{2026-latin-context-matrix-refs}
\end{document}
